\begin{abstract}
针对圆柱形药材烘干中内部干燥滞后、非线性物性与几何收缩共同影响热量和湿度传递的问题，本文基于药材的变物性构建\textbf{径向热湿耦合有限体积模型}，按恒定热物性、变物性和移动边界逐步扩展。根据附件数据计算时变环境与半径曲线，采用全隐式有限体积法求解径向温湿分布，以全截面首次达到判断依据确定干燥终点，并通过同物性对照，定量分析几何收缩的影响。

\textbf{针对问题一：}在固定半径与给定热物性下，建立\textbf{Fourier--Fick 热湿传递基准模型}，以轴线对称条件和\textbf{表面双 Robin 边界}描述内部扩散及环境交换，采用表面加密网格求解预热过程。1800~s 时中心与表面温度分别为{\bfseries\boldmath $33.577~{}^\circ\mathrm C$ 和 $36.786~{}^\circ\mathrm C$}，水分浓度分别为{\bfseries\boldmath $2.550$ 和 $1.510~\mathrm{kg/kg}$}。中心水分仍接近初值，表明预热初期失水主要发生在表层，内部干燥存在明显滞后。

\textbf{针对问题二：}引入热物性对水分、扩散系数对温度与水分的依赖，构建\textbf{变物性非线性热湿耦合模型}，通过 Picard 迭代交替更新场量与物性，反映温度升高与含水量降低对水分扩散的相反作用。3~h 时中心与表面温度分别达到{\bfseries\boldmath $49.849~{}^\circ\mathrm C$ 和 $49.967~{}^\circ\mathrm C$}，水分浓度分别降至{\bfseries\boldmath $1.766$ 和 $1.008~\mathrm{kg/kg}$}。温度趋于均匀而水分仍有明显径向差异，说明水分均匀化滞后于升温过程，温度接近环境值不能作为干燥完成的依据。

\textbf{针对问题三：}将问题二模型扩展为\textbf{基于全截面首达判据的长期干燥模型}。依据末段观测，假设 4~h 后环境保持为 $T_e=50~{}^\circ\mathrm C$、$C_e=0.05~\mathrm{kg/kg}$，采用相邻时间层线性插值，定位全截面水分最大值降至 $0.15~\mathrm{kg/kg}$ 的临界时刻，其后全截面严格达标。得到\textbf{全截面临界达标时间为}{\bfseries\boldmath $57.472$~h}，此时中心与表面水分浓度分别为{\bfseries\boldmath $0.150$ 和 $0.053~\mathrm{kg/kg}$}；\textbf{表面首达时间为}{\bfseries\boldmath $12.552$~h}，早于全截面达标约{\bfseries\boldmath $44.919$~h}。计算表明，药材中心的水分浓度最高、达标最晚；考虑圆柱两端面的传热和失水后，这一结论仍然成立。因此，表面达标并不代表内部已完成干燥，仅依据表面水分判断会导致过早停机。

\textbf{针对问题四：}建立\textbf{收缩移动边界热湿耦合模型}，在均匀材料收缩与干骨架守恒假设下，通过 $\xi=r/R(t)$ 将变化的物理区域映射至固定材料坐标，使材料与网格速度一致，统一处理扩散距离缩短与表面交换变化。得到\textbf{临界干燥时间为}{\bfseries\boldmath $51.088$~h}，\textbf{达标半径为}{\bfseries\boldmath $1.200$~cm}。保持附件物性与环境条件一致，固定初始半径的对照模型需{\bfseries\boldmath $129.845$~h}；在给定物性与收缩轨迹下，\textbf{收缩使干燥时长缩短}{\bfseries\boldmath $60.7\%$}。该对照用于分析收缩效应。

网格与步长加密检验支持数值收敛；\textbf{移动网格验证}中，两种坐标实现的最大水分差为\mbox{\bfseries\boldmath $2.61\times10^{-13}~\mathrm{kg/kg}$}，干燥时间差为\mbox{\bfseries\boldmath $1.07\times10^{-7}~\mathrm{s}$}；\textbf{COMSOL 有限元复核}在六个采样时刻的\textbf{最大水分差为}\mbox{\bfseries\boldmath $4.54\times10^{-5}~\mathrm{kg/kg}$}。在收缩轨迹不变的所设温度与环境水分扰动试验中，\textbf{干燥时长对温度的响应更明显}。模型可确定圆柱药材的干燥终点。

\keywords{药材烘干\quad Fourier--Fick 模型\quad 有限体积法\quad 热湿耦合\quad 移动边界}
\end{abstract}
